\documentclass[a4paper,10pt]{article}
\usepackage[T1]{fontenc}
\usepackage[utf8]{inputenc}
\usepackage[english]{babel}
\usepackage[a4paper,margin=1.5cm]{geometry}
\usepackage{graphicx}
\usepackage{booktabs}
\usepackage{longtable}
\usepackage{float}
\usepackage{pdflscape}
\usepackage{titlesec}
\usepackage{amsmath}
\usepackage{amssymb}
\usepackage{afterpage}
\usepackage{lmodern}
\usepackage{enumitem}
\usepackage{changepage}
\usepackage{multirow}
\usepackage{pifont}
\usepackage{tablefootnote}
\usepackage{threeparttable}
\usepackage{array}
\usepackage[bottom,hang,flushmargin]{footmisc}
\setlength{\skip\footins}{0pt}
\usepackage[toc,page]{appendix}
\usepackage{caption}
\usepackage{stfloats}
\usepackage{microtype}
\usepackage{csquotes}
\usepackage[style=authoryear,backend=biber]{biblatex}
\addbibresource{references.bib}
\usepackage{xcolor}
\usepackage[colorlinks=true,linkcolor=blue,citecolor=blue,urlcolor=blue]{hyperref}
\begin{document}
\clearpage
\onecolumn
\begin{landscape}
\begin{small}
\begingroup
\small
\setlength{\tabcolsep}{2pt}
\renewcommand{\arraystretch}{0.9}
\begin{longtable}{p{0.5cm} p{2cm} p{2.5cm} p{4cm} p{1cm} p{2.5cm} p{5.5cm} p{2.5cm} p{2.5cm} p{1.8cm} p{1.5cm}}
\caption*{Table: Comparative analysis of performance measurement approaches in SRR-AFSCs} \\
\toprule
\textbf{S/N} & \textbf{Author(s)year} & \textbf{Topics focus} & \textbf{Objective(s)}& \textbf{Scope} & \textbf{Methodology} & \textbf{Key observations} & \textbf{Research methods} & \textbf{Limitations} & \textbf{Product} & \textbf{Country} \\
\midrule
\endfirsthead
\multicolumn{11}{c}{\textit{Continued from previous page}} \\
\toprule
\textbf{S/N} & \textbf{Author(s)year} & \textbf{Topics focus} & \textbf{Objective(s)}& \textbf{Scope}& \textbf{Methodology} & \textbf{Key observations} & \textbf{Research methods} & \textbf{Limitations} & \textbf{Product} & \textbf{Country} \\
\midrule
\endhead
\bottomrule
\multicolumn{11}{r}{\textit{Continued on next page}} \\
\endfoot
\bottomrule
\endlastfoot
01 & \textcite{Sorina4} &  Technology roles in enhancing the sustainability of AFSC & develop a comprehensive and systematic analysis of the application of blockchain technologies and related digital solutions in AFSC  & AFSC &  literature review, conceptual approach& convergent integration of emerging digital technologies represent a strategic direction for strengthening resilience, transparency and sustainability by ensuring traceability, fraud prevention, adoption factors, emerging innovations and applications & PRISMA &lack of empirical validation  &NA& NA\\
02 & \textcite{Ayla6} &  sustainability themes  & develop a multi-capital sustainability framework &AFSC & literature review& 116 indicators are identified &PRISMA   & Lack for empirical validation& NA& NA  \\
03 & \textcite{Karel8} &  Applications of the circular economy & identify new trends and gaps&circular AFSC & literature review& the clusters are innovation and sustainability in AFSCs, barriers to sustainable management, strategies to reduce food loss and waste, resource recovery and food security, and conceptual frameworks for waste reduction &bibliometric analysis, SLR, PRISMA   & time scale & NA& NA  \\
04 & \textcite{16Latino} & role of Iot and data analytics in decision-making & provide a road-map for companies to integrate technology for sustainable operations and trust consumer communication&circular AFSC & multi cases approach& data integration facilitates sustainable decision-making in AFSC and facilitate transparent consumer communication&literature review, depth interview, analytic design   & sample & tomatoes, seafood& Italy  \\
05 & \textcite{Reitano22} &  impact of AI in AFSC & explores the transformative potential of AI in the whole AFSC, focusing on its implications for agricultural productivity and consumer behavior & AFSC &literature review& the cluster are: adoption of AI technologies in agriculture, applications of AI in the intermediary management of the supply chain, consumer perceptions and acceptance of AI-powered food innovations, and the role of AI in precision nutrition and personalized health management &PRISMA   & time scale& NA& NA  \\
06 & \textcite{Sajal23} & AI-enabled secure social industrial IOT& analysis the AI-enabled secure social industrial IoT applications in AFSC, addressing key security concerns and efficiency bottlenecks&AFSC & literature review& AI-driven security solutions significantly enhance trust management, anomaly detection, and data privacy in social industrial IoT networks &taxonomy development   & time scale, sample& NA& NA  \\
07 & \textcite{Yohannis25} & contamination in AFSC & synthesis of recent advances in technologies to mitigate contamination in AFSC&AFSC & literature review& integration technologies enables a transformative paradigm shift from reactive monitoring to predictive and preemptive of contamination &critical synthesis   &  data complexity,&cereals, groundnuts, oil seeds, and spices& tropical and subtropical regions  \\
08 & \textcite{Manuela29} &decentralized identity and digital twins for verifiable product identity & identify the impact of decentralized identifiers, digital twins,  smart contracts, and blockchain to create verifiable digital representations of agricultural products and automate trust mechanisms&AFSC & literature review& blockchain technology enhances product identity, traceability, ability to create secure tamper proof digital records, maintains a verifiable history, increases transparency and reduces fraud &PRISMA   &  time scale& NA& NA\\
09 & \textcite{14Guohua} & mode selection and logistics integration strategy &  determine optimal sales mode and logistics strategy for Green Agro-Food Supply Chain (GAFSC) in e-commerce&GAFSC & mathematical modeling& in the reselling mode, when the logistics service fee is low, the retail price and the level of logistics service under the self-operated logistics strategy are higher than under the third-party logistics strategy&game-theoretic analysis  & generalization& NA& China  \\
10 & \textcite{XIAOYA3} &  AFSC network design & design GASCN and minimize total transportation costs &GASCN & mathematical modeling& ESSA outperforms the Genetic Algorithm and Particle Swarm Optimization in fitness and execution time, lower costs and faster convergence &Enhanced Spring Search Algorithm (ESSA), Genetic Algorithm (GA), Particle Swarm Optimization  (PSO) & single-objective (cost) & NA& China \\
11 & \textcite{Marloes18} & management of food waste &  develops a mixed-integer optimization model to support decision-makers in determining an optimal product portfolio and processing configuration&GAFC & mathematical modeling& side-stream valorization is generally well aligned with profit maximization, while it is not always beneficial from an environmental impact perspective &bi-objective mixed-integer programming & generalization& carrot& Netherlands\\
12 & \textcite{Rui1} &  information collaboration in AFSC & constructs a quadripartite game model involving the government, an information service platform, farmers, and agri-food enterprises. Analyzes stakeholders’ strategic interactions and evolutionary pathways of stakeholders while exploring the regulatory effects of key parameters in reward and penalty mechanisms on system convergence&AFSC & hybrid approach &  standalone fiscal subsidies can rapidly initiate the coordination mechanism, but their incentive effects are difficult to sustain over time&game theory, dynamics simulation & lack of accounting for external uncertainties& NA& China \\
13 & \textcite{Firouzeh13} & digital capabilities and their impacts on AFSC resilience & identify the nature and impacts of firms’ digital capabilities on engineering versus socio-ecological resilience &AFSC & mixed methods& AFSC members’ digital capabilities play distinct but interconnected roles in enhancing resilience, with their impact varying across persistence (engineering), adaptation, and transformation (socio-ecological) over different time frames  &qualitative and analytic methods (comprehensive taxonomy) & non-generalization of the results& grains, red meat, dairy, horticulture, seafood, and wine& Australia  \\
14 & \textcite{Juan32} &  Key resilience and sustainability drivers and controllers in AFSC& identify and prioritize the key drivers and enablers of resilience and sustainability in AFSCs&AFSC & mixed approach& the main drivers are "Evolution of Agricultural Systems" , "Water Footprint" . The main enablers are "Redesign and Coordination of Operations" and "Collaboration Across Supply Chain Links" &SLR, AHP, DEMATEL   &lack of empirical validation, limited generalization& NA& NA  \\
15 & \textcite{Yassine10} & blockchain in AFSC& identify barriers to blockchain adoption in the  AFSC and effects for policy guidance&IAFSC & empirical and case study& 12 barriers identified&DEMATEL analysis, expert opinion   & small expert sample& NA& Africa  \\
16 & \textcite{Giovanni9} &   technologies for enhancing transport systems&develop a model that enables
the seamless integration of data from sensors, IoT devices, data management platforms, and distributed ledger technologies (DLT) within a newly designed data space architecture&UAFC, PSC & empirical and case study& platform improve data integrity, provenance, transparency, actionable ML insights from sensor data &design platform   & generalization& cattle and livestock & Italy  \\
17 & \textcite{Conn7} &  AI in AFSC & examine AI-enabled information processing to build resilience in AFSCs via organizational mindfulness and flexibility&PAFSC & empirical and case study& organizational mindfulness and flexibility play key roles in building resilient supply chains, boosting long-term performance, and reducing food waste   & generalization  & NA& Europe   \\
18 & \textcite{Minelle5} &  Sustainable agricultural practices transfer& examine how SUSAPs are implemented and transferred across multi-tier SC using a socio-material, practice‑based lens &UAFC, FPAFC & empirical and case study& animal welfare fully established as a practice; waste management, working conditions, biosecurity remain proto-practices with partial/weak links, first‑tier supplier as boundary spanner &interviews , qualitative methods   &generalization& poultry products & Brazil  \\
19 & \textcite{Abdulaziz11} &  Act on corporate due diligence obligations & assess perceptions and effectiveness of SC using expert perspectives&AFSC & empirical and case studies&large established companies
are perceived to be less affected due to their existing transparency and compliance infrastructure &semi-structured interviews   &  sample& NA& Germany  \\
20 & \textcite{Jayalath15} &   Digitalization of AFSC & develop a digitalized, regulated architecture to transform SCs in developing economies and improve efficiency, transparency, and food security. Evaluate stakeholder interactions&UAFC & empirical and case study&the proposed framework offers benefits to all stakeholders in all supply chain while improving transparency, optimizing resources, leading to improved production planning, minimizing post-harvest wastage, reducing price volatility and increasing food security.   & reasoning approach   & single country focus& seeds, vegetables & Sri Lanka   \\
21 & \textcite{Meisam19} &  Circular economy startups & develop a  circular economy comprehensive typology that identifies their key characteristics and contributions to the CE &UAFC & empirical and case study& three distinct groups identified biotech-driven sustainable producers, digital CE enablers, and circular packaging pioneers&content analysis, hierarchical clustering  & European focus& NA& Europe  \\
22 & \textcite{Rackl21} &  Geographical indication certification & examine the role of Geographical Indication (GI) certification in coordinating small- and medium-sized food suppliers and large-scale retailers in AFSC, where retailers seek to procure high quality goods&AAFSC & empirical and case study& The results suggest a potential role for GI certification in supplier-retailer coordination &survey, SLS , GMM, PSM  & generalization&  Food craft (bakery, butcher, confectioner, miller, brewer, cheese dairy, and winemaker) & Germany \\
23 & \textcite{Vivian2} &Circular economy in AFSC & selection of AFSCs with circular economy potential, through the implementation of quantitative and qualitative prioritization, proposing indicators related to the circular economy&UAFC & multi-criteria analysis& top prioritized products: avocado, grapes, guava, citrus & Shannon and Condorcet index   & regional single region& 44 products (fruits/vegetables) & Colombia  \\
24 & \textcite{Wen31} & Severity of losses and wastes & evaluate the severity of losses and wastes&UAFC & Multicriteria analysis&  the eight severe issues identified  “Inadequate planning and control of overall production and marketing policies” under government policies, “Adverse climate conditions” and “Imbalance in production and marketing” under production site loss, “Inaccurate market demand forecasting” and “Poor inventory management at supermarkets” under supermarket waste, and “Improper storage management of ingredients leading to spoilage” and “Inability to accurately forecast demand due to menu diversity” under restaurant waste. &BWM , expert survey& qualitative weighting with expert subjectivity& NA& Taiwan  \\
25 & \textcite{Sumit24} &  Blockchain implementation challenges &  identify and evaluate the main blockchain challenges affecting food SMEs&AFSC & Multi-criteria analysis& lack of management commitment, negative perception of BCT, and high implementation costs emerged as the primary obstacles to BCT adoption in AFSC&literature review , expert questionnaire , IMF-SWARA , TFBM approach & small expert sample& NA& India    \\
26 & \textcite{Madushan17} &  traditional AFSC &  evaluates the integration of circular economy practices to enhance sustainable production in order to reduce waste in traditional AFSC in developing economies&AFSC & multi criteria analysis& information flows are the most impactful flow in minimizing supply chain waste, followed by material and financial flows&interviews, AHP , ISM/ANP   & single-case & vegetables & Sri Lanka \\
27 & \textcite{Betul12} &  sustainable strategies in AFSC & Assess sustainability AFSC&UAFC & multi-criteria analysis&low R\&D investment is the major barrier while adopting digital technologies emerges is  the most effective solution &SF-PIPRECIA , SF-CODAS , SF-TOPSIS/WISP , SWAM    & single-country case& Tea& Turkey  \\
28 & \textcite{AFFAF1} &  Blockchain in AFSC &  End-to-end blockchain-based AFSC solution with traceability, trading, delivery, reputation, and security&AFSC & mixed approach& Demonstrates traceability, accountability, credibility, and autonomous delivery with vetted gas costs and security analyzes &Smart contracts , IPFS storage, off-chain arbitrator   & scalability& NA& NA \\
29 & \textcite{Federica4} &  role of food waste prevention technologies & investigates the range of the available technologies and the detailed objectives of such technologies for food loss and waste prevention &pre-harvest to post-harvest  & mixed approach&  adoption of different technological options can represent the engine to establish vertical collaborations between the adopter of the technology and another stage in the AFSC, in order to fight food waste and loss with a coordinated supply chain effort &interviews, qualitative taxonomy development & generalization& fruit and vegetable& Italy  \\
30 & \textcite{Samant7} &  Blockchain in AFSC& identify the potential drivers of blockchain technology adoption&AFSC & mixed approach& disintermediation, traceability, price, trust, compliance, and coordination and control can influence the supply chain actors’ adoption-intention decision processes &conjoint analysis   & Generalization& grape wine& India  \\
31 & \textcite{Mishra2} &   disruptions in AFSC & study disruptions in AFSC during COVID-19 and develop capabilities for resilience and operational excellence using SAP-LAP and dynamic capabilities &DAFSC & mixed approach& propositions linking disruptions, capabilities, and resilience, framework showing proactive/reactive strategies, actionable actions for demand/supply/logistics &SAP-LAP framework, interviews, multiple data sources, triangulation & single case,  generalization& perishable products& India   \\
32 & \textcite{Taylor3} & development of lean AFSC & present an initial model of an integrated supply chain based on the application of lean principles&PAFSC& empirical approach& value chain analysis methods combined with lean principles can provide a powerful framework for the analysis and improvement of supply chain activity &value chain analysis (VCA)   & two case studies within a single sector& pork product& UK\\
33 & \textcite{Leat6} &   risk and resilience in AFSC& identify the key risks and challenges involved in developing a resilient agri-food supply system, particularly with regard to primary product supply, and to show how risk management and collaboration among stakeholders can increase chain resilience&PAFSC & Empirical approach&horizontal and vertical collaboration reduced vulnerability, improving market security for producers and supply reliability for retailers &qualitative case study& single specific supply chain analysis& pork products and Pigs& Scotland   \\
34 & \textcite{Rana5} &  Blockchain technology & realize a SLR on Blockchain technology  applications for sustainable AFSC, analyzing benefits and challenge&AFSC & literature review& BCT improves sustainability, traceability, and food safety, but faces challenges like scalability, privacy leakage, high costs, and connectivity &SLR   &  search criteria and keywords excluded & NA& NA    \\
35 & \textcite{Sasi8} &   mitigation in a AFSC& propose a novel heuristic algorithm that revises the ideal plan in the event of a disruption &GFSC & mathematical approach&  proposed heuristic is highly accurate ( gap vs. exact methods) and significantly faster &MILP-H   & study is deterministic and limited to single period& rice and wheat& Global   \\
36 & \textcite{Karki10} &  Closed-Loop supply chain network design optimization  & propose  a sustainable closed-loop fruit supply chain using the Red Deer Algorithm&closed-loop AFSC & mathematical approach&  Red Deer Algorithm (RDA) outperformed other methods (GA, SA, KA, ICA) in solution quality and efficiency &RDA, GA, SA, KA, ICA  &higher computational time (CPU) for large-scale problems&NA \\ 
37 & \textcite{Zheng27} &  impact of environmental regulations and digital empowerment & Understand environmental regulation and digital empowerment on green AFSC under uncertainty using a tripartite stochastic evolutionary game&GAFSC & mixed approach& Existence of multiple equilibria; digital tech and regulation synergistically drive green transformation &Stochastic evolutionary game theory , Gaussian white noise  & generalization& chicken& China  \\
38 & \textcite{kafi2} &  AFSC Research Directions & Identification of future research directions & AFSC & literature review& 7 main  research clusters identified: AFSC system; AFSC management; AFSC industry; AFSC risk factors; AFSC information; AFSC advancement; AFSC risk management & bibliometric analysis &recent publications not reviewed (until 2021)  &NA& NA  \\
39 & \textcite{sharma7} &role of I4Ts in enhancing the resilience of AFSCs & analyze resilience enablers and assess industry 4.0 technologies (I4Ts)  & AFSC & literature review, integrated methodology, interviews & the most influential resilient enablers: real-time information sharing, enhanced product traceability, improved risk management, and planning and network design; I4Ts enhance the resilient enablers: big data analytics, Internet of things, and cloud computing  &  GINA: identify the most influential resilience enablers; FLQOWA: assess and prioritize I4Ts to enhance resilient enablers& no verification of causality of factors, absence of data-cleaning, non-generalization of the results & NA& United Arab Emirates\\
40 & \textcite{khan8} &  role of intermediaries in the sustainability of AFSC& assemble all the stakeholders of an AFSC into a coordinated platform & MTAFSC & single case study methodology & this integration has contributed to increase actors' trust, collaboration at different levels   & qualitative case study method based on semi-structured interviews & lack of evaluation of the economic impact of the sustainable practices on costs in  MTAFSC& rice& Pakistan\\
41 & \textcite{percin10}  & integration of industry 4.0, sustainability, and circular economy (I4.0-S-CE) into AFSC&identify key I4.0-S-CE drivers and the prioritized order of importance of agricultural sustainable development goals  & AFSC & literature review, expert opinions, case study& the most important sub-drivers: resource efficiency, economic and social sustainability, and the achievement of standards and SDGs . The effective use of resources is considered the most important key to meeting the agricultural needs of countries. Sustainable consumption and production are the most crucial targets. Then, hunger, clean energy, and life on land will guide countries towards the goals of zero hunger, renewable and clean energy, and ecosystem protection. Furthermore,  other agricultural SDGs are essential, such as eradication of poverty, improved quality of life, protection of water resources and mitigation of climate change  & hybrid Fuzzy MCDM Method (FBWM, FTOPSIS, FARAS, FSAW)& not taken into account the peculiarity of developing countries (insufficiency of technological infrastructure, lack of cooperation between stakeholders, and limited awareness of sustainability, circularity and regulatory) & Oils, dry food, bakery products,canned products,drinks & developing countries (Türkiye) \\
42 & \textcite{veloso11} &  assessing circular economy in AFSC&identify strategies that promote EC to various stages of the AFSC, identify a classification of sustainability KPIs & AFSC &case research methodology (analysis of sustainability reports)  & 31 indicators identified for measuring, supporting, or controlling CE strategies+6 new standards proposed: avoided food waste; food waste identification and management; reverse logistics; circular packaging design; and regenerative farming& expert focus group, and practitioners’ interviews & research not exhaustive, limited sample & NA & Portugal  \\
43 & \textcite{obayi12} &Intersection of industry 4.0 technologies and human–machine value&explore how system and business-level interoperability at the’ motivations & Agri-food digital ecosystem & qualitative research approach & system-level is driven by defensive strategies for compliance and standardization, while the interoperability at the system level relies on offensive strategies for flexibility and innovation& case study, qualitative analysis (interviews, field observations and focus group discussions) & generalization (focused on one digital ecosystem in India)  & NA& India  \\
44 &  \textcite{zhao14} & Adopting I4.0 technologies in AFSCs&identify barriers to the adoption of I4.0 in AFSCs. & AFSC & fuzzy analytic hierarchy process & the key challenges to applying I4.0 are the increased cost of terminal logistics, acquisition of intelligent agricultural equipment subsidies based on Guanxi, low compatibility of I4.0 technologies with existing agricultural equipment, and problems with the government subsidy model& literature review , consultation of experts and Multi-criteria tool & Generalization, limited scope of literature and few expert input.& NA& china\\
45 & \textcite{lipinska16} & Protection of the position of an agricultural producer in AFSC& determining how the legal instruments provide protection for agricultural producers against unfair trade practices to strengthen the position and increase the competitiveness of the agricultural producer & PAPSC & qualitative and legal methodology & strengthening the position of farmers in the food supply chain is one of the key objectives of the Common Agricultural Policy (CAP). It can be achieved through the development of cooperation among the actors of the chain (agricultural producers)& legal analysis, comparative study, interpretation & study focused on the analysis of legal texts, without direct treatment of the concrete effects on producers & NA& NA \\
46 & \textcite{Tebaldi17} &quantitatively assessing the resilience level&propose an analytic quantitative model to quantify the resilience level of an AF system & AFSC & analytic model, case study & 23 qualitative and quantitative indicators identified for a resilient AFSC& literature review, descriptive statistics method   & Employing in-depth interviews, study geographically limited & coffee& Ethiopia \\
47 & \textcite{Habib18} & role of critical factors in the durability of AFSC&Identify drivers, enablers, and barriers to adopting sustainable practices & AFSC & literature review(identifying factors) and empirical survey(ranking factors) & The key critical factors : economic and productivity improvement, cost effectiveness and improvement in overall performance, and difficulty in mindset and cultural changes & descriptive statistics method (analyze data+ranking factors) & Employing in-depth interviews, study geographically limited & coffee& Ethiopia\\
48 &\textcite{ElHajji19} & Optimization AFSC  using platform based on blockchain&propose a web platform "Hyperledger Fabric" for dynamic real time visibility of a package’s movement & AFSC & platform-based design and evaluation &  real time visibility of a package’s movement and security of the AFSC & cryptographic algorithms (MSP and CA), evaluation metrics, benchmarking results & the lack of predictive insights which can be resolved by on integrating  IoT  and advanced analytics& NA& NA \\
49 & \textcite{zhao20}& Obstacles to adopting Industry 4.0 for a sustainable AFSC&identify barriers and prioritize them & AFSC &SLR,  structured interview+ group-based fuzzy analytic hierarchy process (GFAHP)& The three main barriers are I4.0 technologies that increase the cost of terminal logistics, the acquisition of subsidies for smart agricultural equipment and the incompatibility of technologies with existing agricultural equipment&   literature review, survey &  generalization, limit to papers published in high-quality journals, three experts& NA& Chine\\
50 & \textcite{ricciolini22} &Sustainability of European AFSC& measure the sustainability  level (by countries and sectors: agriculture, industry, distribution and consumption) & AFSC & multi-criteria analysis &Italy ranked highest in AFSC sustainability, followed by Sweden, Austria, Spain, France, Germany, Czech Republic, Portugal, and Slovakia. The most significant sustainability issues were observed in the consumer sector& Multiple Reference Point Partially Compensatory Indicator MRP‑PCI & generalization& NA& European Union  \\
51 & \textcite{Sharma23} & selection of performance indicators for sustainable supplier evaluation& propose a group consensus decision making model based on an iterative algorithm to generate weights, select suppliers who can embed factories ’s sustainability requirements into their SC operations& AFSC & multi-criteria decision-making & the top ranking indicator for sustainable supplier evaluation research is ‘Food Safety’ then ‘Payment to Sub-suppliers’& Delphi method, GAHP, Kendall’s concordance coefficient, case study &  ranking of performance indicators may change according to the decision-making environment &flour milling & India \\
52 & \textcite{belhadi24} & Management of AFSC uncertainties following geopolitical disruptions& study the role of digital dynamic capabilities in enhancing SCRes in  AFSCs impacted by geopolitical disruptions& AFSC & multi-case study approach & strong relationship between the dynamic capabilities of digital supply chain technologies and SCRes at both organizational and inter-organizational levels, proactive and reactive SCRes are built on digital capabilities & qualitative, exploratory research, multi-case study design, interpretive analysis, exploratory interviews, semi-structured interviews, cross-case analysis& generalization (14 firms)& NA& Africa  \\
53 & \textcite{Rackl26} &  Coordination in AFSC through geographical certification&identify the impact of geographical indication certification for vertical coordination & AAFSC&Interviews, empirical study &GI certification enhances the coordination of procurement contracts by improving the ability of the retailer to identify efficient suppliers and encouraging the provision of high-quality goods& econometric model & non-generic results  depending on the product type and country & food craft & Germany  \\
54 & \textcite{Chandrasiri29} &Adoption of bi-modal transport for reducing food miles & reduce transport distances (food miles and and empty haul distance), transport cost and emission footprints by modifying the volume and direction of the product flow & DAFSC & analysis model, case study & bi-modal transport allow the reduction of food miles by 32\%, transportation costs by 36\%, contributions to the global warming potential by 35\%, and empty truck hauls by 38\% & mixed-integer linear programming model & single-objective linear programming, generalized constraint (different categories of vegetables, storage condition, shelf-life time  of vegetables,etc) & vegetables& Sri Lanka \\
55 & \textcite{GamezAlban30} & vehicle routing problem (problem of the pickup round: first mile): Heterogeneous Fleet Routing Problem (HFRP) & minimize total logistics costs of visiting fields during a pickup round using multiple vehicles. &PHL& analytic model+case study  & organizing the first mile with a cooled vehicle for pick-up rounds reduces logistics costs by up to 40\% \% compared to other practices& mixed-integer linear model& Single-objective linear programming (e.g. lack of study of maximization of product freshness at the time of sale ) &peaches and nectarines&  Greece \\
56 & \textcite{Aghajani31} &optimization AFSC & maximize the net value of the whole supply chain, maximize the job opportunity, minimize the greenhouse gas emission (minimize overall cost at various stages, maximize the revenue and employment generated and reduce the environmental impacts, maximize job opportunities in the chain by maximizing the man hours needed while reducing the greenhouse gas emission, propose a pricing model to estimate the price in various customer zones) & AFSC & literature review,  mathematical modeling, sensitive analysis, robust possibilistic programming& temperature exert the highest impact, whereas land use shows the lowest& multi-objective integrated mathematical model &  model limited to a single case study and expert-based assumptions. & stevia& Iran \\
57 & \textcite{Cahyadi32}  & integration of circular economy (CE) principles into the AFSC& identify  existing knowledge, research gaps & AFSC & systematic literature review& CE enhances sustainability and efficiency within the agri-food sector and generates economic and environmental benefits on a regional and global scale &reviews and Meta-Analyses (PRISMA)& rely on AI for abstract Screening, text review and quality assessment, the selection criteria (most cited articles)& NA& NA \\
58 & \textcite{Arimany33}  &  causal relationship between resilience and sustainability & identify the sustainability variables that help the resilience of companies & RSC &qualitative comparative analysis methodology, case study&  the resilience  of supply chains is strongly linked to strategies to improve environmental, social and governance sustainability (pollutant emissions, safety and health of workers, diversity in management bodies and employees, reduction in energy consumption,  selection of suppliers with social criteria and  fight against corruption)  & investigation &  sample of eight distribution companies&beverages, frozen foods, dairy, dry foods and grocery, meat and charcuterie, pet food, home care, personal care & Spain\\
59 & \textcite{Sakthivel34}  & Blockchain enhances trust and traceability & propose a blockchain architecture & CAFSC & experimental methodology& blockchain enhance managing a massive number of data and transactions, avoiding fraud, securing records, strengthening trust & architecture, deployment, evaluation &lack of traceability  & NA& NA \\
60 & \textcite{Charlebois35}  &digital traceability adoption and its impact in AFSC & provide an overview of countries’ efforts in digital food traceability and propose actionable recommendations and insights for effective deployment of digital traceability systems& AFSC & comparative analysis, multidimensional analytical framework&  OECD countries increasingly adopt digital traceability especially blockchain and QR codes enhancing food safety and transparency, while highlighting disparities, regulatory gaps, and the need for standardized framework & case studies and empirical study( quantitative analysis, regression analysis) & incomplete data, generalization&NA &38 Organization for Economic Co-operation and Development (OECD) member countries \\
61 & \textcite{Sayma36}  & reason of price hike in AFSC utilizing private Ethereum blockchain network & propose a system that enables tracing the origin of corruption in the supply chain and identifying the source of price increases & AFSC & conceptual framework, system design, implementation & the proposed system ensures a high level of transparency, as every transaction recorded on the blockchain is visible to all network actors, thereby promoting openness and reducing information asymmetry. it provides strong traceability, allowing each transaction to be tracked back to its specific Ethereum address. This feature facilitates Data Protection , verification, Fraud Prevention,Security ,Transaction speed and Cost Efficiency across the entire supply chain.& private Ethereum blockchain, Proof of Authority (PoA) approach, InterPlanetary File System (IPFS), dApp (decentralized Application)&  simple and no complex AFSC  & vegetables (potato and eggplant)& Bangladesh  \\
62 & \textcite{Wang37}  & Impact of digital technology in the sustainability of CBAFSC & identify the research hotspots and the potential research directions & CBAFSC & bibliometric analysis& digital technology enables sustainability and food safety in CBAFSC & bibliometric analysis with  visual knowledge map & sample literature, interpretation with “Citespace ”, without comparing and validating the results with other literature analysis software (eg. VOSviewer  , Histcite) &  NA& NA\\
63 & \textcite{Elbhilat38}  &Impact of AI on AFSC performance &examine the effect of AI based technologies on reducing wastes and minimizing AFSC costs, hence increasing organization profitability & AFSC & conceptual model, Harman statistical test, Partial Least Square Structural Equation Modeling (PLS-SEM), confirmatory factor analysis (CFA), quantitative approach, common method bias (CMB), hypothesis test (survey), dynamic capability view (DCV) and organizational information-processing theories& distribution network efficiency ,AFSC efficiency, and organization performance are  directly and positively affected by AI integration. Efficient distribution logistics in the agri-food sector reduces product waste, ensures timely delivery of perishable goods, and improves response to market changes. &Partial Least Square-Structural Equation Modeling (PLS-SEM) approach& study based on survey study, generalizability & NA& Morocco \\
64 & \textcite{Camel39}  & integration of BCT with Green Product Platforming (GPP) in the agri-food industry & analyze the impact of BCT-enabled Green Product Platforming on reducing the carbon footprint and enhancing economic performance & AFSC& quantitative surveys, qualitative interviews, data analysis & the importance of aligning GPP initiatives with stakeholder expectations and leveraging BCT to navigate and capitalize on stakeholder pressures for improved carbon and economic performance& partial least squares structural -equation modeling (PLS-SEM). & sample, generalization  & NA& NA \\
65 & \textcite{Hamidoglu40}  & government-backed agri-food supply chain model (GBASM) between  farmers and agricultural enterprises &  assess the impact of the GBASM on the decision-making processes of farmers and agricultural enterprises (supplier, manufacturers/factories, distributors, supermarkets/wholesaler/ retailer) & GBAFSC & game-theoretical approach + case study& the proposed model allows to achieve affordable market prices among all relevant actors and establish efficient pricing mechanisms within their market environments &Nash equilibrium +numerical experiments, grey wolf optimizer algorithm& Increased complexity (increasing the number of actors  makes the Nash equilibrium more difficult to establish)  & tomato& Argentina \\
66 & \textcite{Liu43}  &Impact of government tax subsidies on Blockchain technology investment strategies in AFSC & analyzes tax subsidies for blockchain adoption and updates the demand function based on consumer trust & GAFSC & Game-theoretic modeling&If blockchain investment costs are within a certain range, tax subsidies yield higher stakeholder benefits. However, full subsidies do not always raise all actors’ incomes, as this depends on subsidy rates and the blockchain technology investment cost & Mathematical modeling, reverse analysis (master-slave game), Numerical simulation & not distinction between online and offline channel& cherries&China \\
67 & \textcite{Guidani46}  & deployment of a digital twin &Virtualize AFSC processes to estimate economic, logistical, environmental, safety, and nutritional indicators for every food order. & AFSC&deployment & the digital twin provides process visibility and helps practitioners make operational and tactical decisions through real-time monitoring and multidimensional post-performance analysis. It also enables consumers, through their informed choices, to exert strategic pressure on the food industry to ensure sustainability.& Relational modeling (SQL) +simulation +OO paradigm +Monte Carlo method  & limited simulation analysis (performed on four virtual scenarios) & nine different horticultural products& Italia  \\
68 & \textcite{Toader49}  & Blockchain technology in AFSC&analyze the factors exerting a positive effect on blockchain adoption within AFSC & AFSC & empirical study &the significant influence of performance expectancy , effort expectancy , and perceived trust on usage intention, while also revealing the positive impact of Organizational Blockchain Readiness (OBR) on expected Usage Behavior (UB) &Unified Theory of Acceptance and Use of Technology(UTAUT)) , Partial Least Squares Structural Equation Modeling (PLS-SEM), investigations (questionnaire), Common Method Bias (CMB)& Sample, generalizing, non-exhaustive study factors & livestock and fresh products&  European countries (Czech Republic, France, Germany, Italy, Poland, Romania, Spain, and Switzerland)\\
69 & \textcite{Marin50}  & impact of AFSCs in biodiversity of in peri-urban agricultural Areas& analyze the role of CFSc in conserving and promoting local agro-biodiversity  & CFSC  & case study approach&distribution plays a crucial role in local food systems, not only involves the trade of food but also influences consumption patterns, societal values, and the overall landscape & semi-structured interviews, inventory& local scope, limited sample & fruits and vegetables& Catalonia (Spain) \\
70 & \textcite{Chen51}  & influence of critical factors on the success of knowledge mobilization& identify  the critical success factors and the correlation between them & IAFSC & quantitative analytical methodology& 8 of the identified criteria are strongly correlated with the success of knowledge mobilization (collaboration, structure of the supply network, power, technology, trust, commitment, training/education, time, cost, culture, and continuous improvement), 3 of them disappeared ( continuous improvement, time and cost)  & questionnaire surveys, multiple regression analysis  & rank the factors using more scientific methods, generalization, model based on literature& NA& NA \\
71 & \textcite{Cossu52} & Eco-labelling and sustainability AFSC & identify the main impact categories of the environmental footprint of the product life cycle. & AFSC &Empirical Methodology(Case studies+ analysisis & The score reflects an intermediate environmental performance +the role of eco-labeling in encouraging more sustainable practices in the sheep dairy phase +the crucial role of sheep farming systems for improving the environmental footprint . & PEF method+ experimental MGI (Made Green in Italy) application+ ACV& sample, generalization&hard sheep milk cheese& Italy \\
72 & \textcite{stevens54} & impact of diversification firms in resilience AFSC &  develop revenue-based measures of firms’ vertical (across supply-chain segment) and horizontal (within supply-chain segment) diversification.  Compare post pandemic outcomes of more diversified and less diversified firms  &AFSC & empirical approach & vertical diversification reduces firms’ resilience, whereas horizontal diversification increases firms’ resilience, horizontal diversification of small and medium sized firms increases resilience in the AFSC &inverse probability weighting (IPW), CBGPS and NPCBGPS techniques, survey& sample& NA& U.S (California, Florida, and the two-state region of Minnesota–Wisconsin)\\
73 & \textcite{Dhita55} &risk management and sustainability in the AFSC & identify the development of research topics and the future research opportunities  &AFSC & bibliometric analysis& increasing trend in publications on risk management, sustainability, and AFSC over the past six years, future research shoud focus to the impact of emerging technologies in resilience and sustainability supply chain &bibliometric analysis & single Scopus, based database& NA& NA \\
74 & \textcite{chu56} & blockchain-enabled vertical coordination in AFSC &develop a blockchain-based framework for improving vertical coordination in AFSC & IAFSC& literature review, framework, case study &blockchain technology provides benefits for each supply chain actor in terms of operational effectiveness, cost and time savings, and improved human resource&conceptual framework analysis & No real-world implementation or empirical validation of the proposed framework& cashew nut& Vietnam \\
75 & \textcite{Boubakeur57} &blockchain for Secure and Traceable AFSC & propose lightweight blockchain-based signature mechanism for rapid detection of food fraud&AFSC & Semantic-based Approach& \~7.5 speed improvement over iterative search algorithms and maintaining high transaction traceability with significantly reducing storage costs &systematic analysis, semantic approach& no privacy-preserving technologies, no advanced detection of unknown attacks & NA& Algeria  \\
76 & \textcite{Bosung58} &improving product safety & examine  direct investment in supplier development and the adoption of blockchain technology for enhanced traceability and transparency, the impact of an increased inspection frequency on food safety, verification the current government interventions (increasing inspection frequency and raising quality standards) &IAFSC & qualitative approach and game theory analysis& both direct investment and blockchain technology adoption can independently promote agrifood product safety, blockchain technology adoption is effective when the retailer secures a low unit margin (i.e., a high wholesale price)&intensive interviews, Nash equilibrium & lack of demand uncertainty of the retailer and considering competition with other retailers& NA& Korea  \\
77 &\textcite{Barbieri59} & the assess the environmental performance of the logistics distribution processes& identify key elements for evaluating the environmental impact of distribution processes & AFSC &empirical-analytical simulation approach& intermodal transport as more environmentally sustainable than full road transport. Potential CO2 emission savings of up to 50\% through improved vehicle efficiency and optimized vehicle utilization. Additional savings of up to 30\% through increased use of intermodal transportation. &semi-structured interviews, algorithmic model, business case  & limited scope  & NA& EU \\
78 & \textcite{otter60} & digital transparency in AFSC networks & identify the primary and secondary stakeholders in the AFSCNs of the digital era & AFSC & empirical study& 11 secondary net chain stakeholders must be integrated in information AFSC network & qualitative semi-structured interviews with technology providers, structured  analysis & lack of assessment of the importance and interests of stakeholders & NA& EU \\
79 & \textcite{Yin61} &relationships between AI, AFSC management practices, resilience, and performance & study the impact of AI on  AFSC  resilience, and performance &PAFSC & analytical approach, empirical study, statistical analysis& the top organizational factors play a crucial role in ASC resilience through management support (TMS), supplier integration (SI), internal management (IM), customer integration (CI), and supply chain collaboration (SCC). The  AI moderates these relationships, AI improves internal management practices due to advanced analytical capabilities for effective decision-making (prediction, and improvement  of communication, efficiency  and coordination )& partial least squares structural equation modeling (PLS-SEM), survey , Harman factor, analysis of reliability metrics& non generalizable& NA& China  \\
80 & \textcite{Moreno62} &correlation between coordination mechanisms  and sustainability performance&evaluate the role of coordination and transaction costs in the sustainability performance SC, examines the interrelations between transactional characteristics, relational and formal coordination, and sustainability performance in AFSC &AFSC & exploratory and quantitative study, statistical analysis & relational governance improve economic and social sustainability, enhance indirectly  ecological performance via resource efficiency+, formal coordination and effective transaction cost management enhance collaboration and supply chain viability&survey, structural equation modeling, confirmatory factor analysis, exploratory and confirmatory factor analyses, structural equation modeling, validity and reliability testing, structural equation model (SEM), Promax Rotation& limited number of variables to measure sustainability performance& blackberry& Ecuador  \\
81 & \textcite{Paredes63} & sustainability and resilience in AFSC & identify the relation between resilience and sustainability&AFSC &literature review& resilience and sustainability as two cross cutting and highly interrelated concepts in supply chain management &  systematic reviews and Meta-Analyses (PRISMA protocol) & content analysis sample of papers, lack of a robust relationship& NA& NA\\
82 & \textcite{Loy64} &environmental impact reporting and  Software solutions in AFSC & determine the importance of developing and implementing advanced supply chain management software to address the need for improved sustainability practices in AFSC, identify the SCM software is being used for enhance sustainability,  determine the impact of software skills on sustainability reporting and governance &AFSC & qualitative research approach& the adoption of software allows for a significant reduction in data collection costs, but human intervention remains essential to manage technological failures &semi-structured, qualitative interviews & limited sample size& NA& United States  \\
83 & \textcite{Israel66} & farm economic viability & identify the interaction between agricultural supply chain capabilities (AGSC), dynamic pricing (DY PG) and subsidy schemes (SBSM) and their collective influence on farm economic viability &LAFSC & empirical approach& smallholder farmer face various challenges that impact their agricultural supply chain capabilities, generating significant increases in food market prices (eg. increased costs related to transportation, storage, and losses) &survey, structural equation modelling (SEM) , structural equation modelling (SEM), descriptive statistics and correlation analysis, pearson correlation, moderated mediation analysis & generalization& staple food products (millet, wheat, barley, maize, rice, and sorghum)) & Tanzania (Dodoma, Morogoro, Mbeya, Songwe, and Tabora) \\
84 & \textcite{Hasan70} &IOT on enhancing transparency and efficiency  of AFSC &define the influence of IoT on Physical and Information Flow(PHF) and the impact of an IoT-enhanced PHF on chain transparency, this helps promote better decision-making, efficiency and consumer confidence &AFSC & experimental design approach, hypothesis Testing and Path Analysis& significant and positive impact of IoT technologies on the PHF and on the transparency of AFSC &sample (probability simple random sampling), exploratory factor analysis,  confirmatory factor analysis (CFA), structural equation modeling (SEM) & generalization, sample size& NA&  Bangladesh \\
85 & \textcite{Gonzalez71} &key issues and trends of blockchain in AFSC & identify future research and gaps in this topic&AFSC & bibliometric analysis& the research hotspots are related to the integration of blockchain technology in the AFSC for traceability, coordination, transparency and Enhancement of food safety &co-occurrence analysis, descriptive analysis, bibliometric analysis & limited database, time scale& NA& NA  \\
86 & \textcite{Joshi75} &role of Industry 4.0 in enhancement the sustainability of AFSC during the health crisis & assess the impact of digital technology 4.0 on sustainability of AFSCs &AFSC & quantitative research&  industry 4.0 technologies have major impact on the viability of sustainability of AFSC.  Organizational flexibility determines how Industry 4.0 impacts AFSC sustainability&IPT theory, structural equation modelling ,PLS-SEM, statistical analysis (measurement model, structural model, hypothesis testing)& limited sample (including companies in transition)& food grains &  India \\
87 & \textcite{Gholian76} & design a sustainable closed-loop AFSC network & identify the optimal number and location of facilities, the flows (i.e., the number of goods) to be transported based on facilities’ capacity, and the inventory level of products, minimizes the total cost, environmental pollution, and maximizes job opportunities&CLAFSC & depth metaheuristic  & optimizers can solve all sizes of problems &non-dominated sorting genetic algorithm-II (NSGA-II), interval plots and the Friedman statistical test, six metaheuristics and three hybrid metaheuristics are utilized, Taguchi method, sensitivity analysis& deterministic data, static model& coconut industries& Mexico  \\
88 & \textcite{Zhao77} &resilience strategies & identify resilience capability factors used at AFSC levels in the face of the  crisis&AFSC & empirical approach&  frequently discussed resilience capabilities, such as collaboration, redundancy, flexibility, leadership, and innovation, were implemented across the preparedness, response and recovery, and adaption phases; however, successful AFSC recovery also depends on each country's cultural value orientations &middle-range theory, thematic and comparative analyses, middle-range
theory (MRT)  & limited geographic scope& NA& China and Spain  \\
89 & \textcite{Balezentis78} &effects of the disruptive events on the viability of AFSC & assess of  AFSC viability during the disruptive external effects (effects of COVID-19 and the Ukraine war)&APC & probabilistic model& most of the negative factors arise from increased energy consumption in the agri-food sector. Positive effects are linked to the increase in production, which should  associated with  the ability of Lithuanian agri-food producers to maintain production activities uninterrupted&complex and sophisticated expert evaluation processing technique, power ordered weighted average (POWA) operator, Monte Carlo simulation   & limited scope of actors& cereals, horticulture, fruit and berries, dairy, poultry, beef cattle, and pigs& Lithuania  \\
90 & \textcite{Trevisan79} &digital technologies for prevention and reduction of food loss and waste & analysis the state-of-the-art of adoption of each digital technologies for preventing and reducing FLW&AFSC & systematic literature review and research agenda& digital technologies are more frequently adopted in the earlier stages of the supply chain, IoT and digital platforms are the most frequently cited technologies, IoT is diffused throughout the entire supply chain&SLR   & generalization & NA& NA  \\
91 & \textcite{Ana81} &improving the robustness in AFSC &  design a a system dynamics (SD) model to analyze, measure and improve the robustness of an configuration and its planting planning to disruptions in demand, supply, transport between SC nodes, and the operability of SC nodes& PASC & mathematical formulation and simulation of scenarios& in the event of a disruption (logistics bottlenecks are the most critical), demand disruptions have immediate impacts on satisfaction and revenue, adaptation strategies (opening alternative nodes, adjusting workforce, reallocating logistics flows&known behavior replication test and the extreme conditions test, case study &generalization& vegetables& Spain  \\
92 & \textcite{Zhao82} &links between risks and resilience capabilities& identify key risks and resilience capabilities in AFSCs, the relationships, correlations and causalities between them, and research gaps and future research directions in the field&AFSC & systematic literature review (SLR), comprehensive analysis& identify eight types of AFSC risk and seven types of AFSC resilience capability &content analysis& lack of an empirical study& multi-product& multi-country  \\
93 & \textcite{Guoqing83} &modelling enablers & identify the  resilience capability factors used to build AFSCRes, identify capability factors interrelated, identify resilience capability factors key to building AFSCRes  &CBAFSC &cross-impact matrix multiplication applied to classification (MICMAC) analysis, cross-country comparative analysis& contractual restraints regulating farmers’ opportunistic behavior and regular interactions are key factors for building AFSCRes in France and Argentina, respectively &interviews, thematic analysis, total interpretive structural modelling (TISM)   & validity of the empirical findings& organic Vegetables, tomatoes, cereals, seeds, agricultural inputs &Argentina and France  \\
94 & \textcite{Kersten84} &traceability in AFSC & investigate the traceability and its technologies&CBAFSC & bibliometric and content analysis( systematic literature review SLR)& traceability is essential at all levels of the supply chain, contributing to the management of food waste and sustainable practices. Various technologies are presented and a framework was developed to illustrate the application of these technologies at different stages of the supply chain and their relationships with the dimensions of ReSOLVE &systematic and structured process, PRISMA (Preferred Reporting Items for Systematic Reviews and Meta-Analyses)  & uses only the ReSOLVE framework to analyze technologies& NA& NA  \\
95 & \textcite{Körmendiová85} &Vertical linkages in AFSC &identify the factors that influence the choice of distribution channel of family farms  &FFFSC & statistical analysis& the choice of the distribution channel depends on the size of the  company and the type of production. Businesses with less than 10 employees prefer direct distribution . Direct distribution channels are used primarily in the supply chain of animal products contrary to the indirect distribution chain specified mainly in the supply chain of plant products &questionnaire survey, statistical analysis, chi-square test & sample& NA& Slovakia  \\
96 & \textcite{Wang86} &cross-border E-Commerce platforms in AFSC & identify and classify the key facilitating factors influencing the sustainable development of e-commerce platform & CBAFSC & systemic approach& the five most fundamental factors are simplification of customs clearance procedures, digital empowerment, consumer demand forecasting, knowledge innovation linkage, and policy environment optimization  & Grounded Theory and DEMATELISM- MICMAC  & non-taking of heterogeneity of platforms& NA& china  \\
97 & \textcite{Krstić87} &e‑traceability drivers & identify factors that motivate companies to implement electronic systems for tracking (e‑traceability) and monitoring products&AFSC & multi-criteria decision-making (MCDM)& the most important drivers are supply chain efficiency, technology development and sustainability &"Factor Relationship" (FARE) and “Axial Distance based Aggregated Measurement” (ADAM) methods   & sample of 27 experts& NA& NA  \\
98 & \textcite{Gholian88} &design of AFSC network& develop a mathematic model to manage the total net profit while monitoring CO2 emissions and the satisfaction of customers within the network&CLAFSC & math& optimizers can solve all sizes of problems. However, MOGWO is better suited for addressing small-size problems, whereas MOHHSA is highly effective for tackling problems of medium and large sizes &mixed-integer linear programming model, conventional and hybrid metaheuristics (four multi-objective optimizers and three hybrid algorithms are)  & lack of reliable database, environment uncertain not modeled& multi product soybean& Iran  \\
99 & \textcite{Rahbari89} &strategic analysis for AFSC under crisis & develop a model of 
food supply chain under uncertain conditions in order to analyze it strategically during the pandemic, optimize production, storage, transportation, resources and exports in the supply chain&SSFSC & MILP& the export of canned food to neighboring countries with economic justification” was the best strategy for the company under study during the COVID-19 pandemic. this strategy reduce b supply chain costs and
increase  the human resources employed. Finally, the utilization of available vehicle capacity 
 and the utilization of available production are optimal  when using this strategy &SWOT method, PORTER and
PESTEL, External Factor Evaluation (EFE) Matrix, the Internal Factor Evaluation (IFE) Matrix+ MILP +TOPSIS technique+ SEM, & generalization, one company in the study& canned food & Iran  \\
100 & \textcite{Say90} &impacts of the pandemic on the AFSC& identify the challenges of AFSC, effective alternative strategies to be implemented by the government &AFSC & desk review& working under strict regulations came with a huge price paid for agrifood industries. This pandemic has affected the  product availability and accessibility. Therefore, the government should be ready with the roadmap and enforce the measures to manage emergencies without disrupting the AFSC &literature review & limited geographic scope and time scale& NA& Malaysia  \\
101 & \textcite{Dong91} &impact of digitalization on AFSC  & investigates the relationships between digitalization and supply chain elements, particularly integration, communication, operation, and distribution, and their effects on corporate profitability and sustainability&IOAFSC & quantitative approach& synergistic effect of digital advancements leads to increased business profitability and sustainability &empirical investigation, PLS-SEM, statistical approach& study based only on a quantitative design& NA& Pakistan  \\
102 & \textcite{Belamkar92} &optimization of AFSC networking problem &develop a multi-objective model (multi-objective MILP model). Aims:  maximize transportation profit, reducing transportation cost, reducing carbon dioxide emissions &PDL &multi-objective optimization problem (MOOP)& proposed model help to achieve balanced solutions between profit, cost and emissions & Weighted sum method (WSM), Fuzzy goal programming (FGP), Interactive Fuzzy satisfying technique (IFST), Global criterion method (GCM), and Weighted Tchebychef metrics programming (WTMP), study case  & generalization, single period& apple& India \\
103 & \textcite{Azab93} &modeling of sustainable AFSC with product perishability and environmental considerations & propose a model to maximize total supply chain profit and minimize total CO2 emissions&processing AFSC & mathematical approach& selection of plots to be planted, selection of the harvesting period, calculation of the transported quantity, determination of the processed quantity of products, estimation of the total stored quantity, determination of the satisfied demand and lost sales&multi-objective mixed-integer linear programming model MILP, study case, Pareto frontier& the model only optimizes economic and environmental dimensions& sugar beet& Netherlands  \\
104 & \textcite{Wang94} &visual knowledge map analysis & identify research hotspots and development
trends, clarify the historical context, and explore future trends  &CBAFSC & bibliometric analysis& since 2017, research in this field has shown a strong development trend, the resilience of the global food supply system have become the main lines of research in this field, research hotspots included cross-border agricultural product safety, cross-border agricultural supply chain systematization, and technology integration&PRISMA  & sample literature from 2007 to 2022& NA& NA  \\
105 & \textcite{Karakoc95} &Structural choke points & identify the spatially  structural choke points &AFSC & complex network approach& the 7 structural choke points identified are generally consistent through time, particularly for processed food commodities& network statistics   & perimeter& live animal and fish, cereal grains, agricultural products, animal feed, eggs, honey and other products of animal origin, meat, poultry, fish, seafood and their preparations, milled grain products and preparations, and bakery products, other prepared foodstuffs, fats and oils& United States  \\
106 & \textcite{Yang98} & the initial real information sharing in the AFSC  & develops an evolutionary game model of AFSC participants +the impacts of the key parameters on the dynamic evolution process of participants&AFSC & modeling, simulation& the greater the risk coefficient of using blockchain technology, the higher the information cost of sharing initial real information, When the reward incentive coefficient,
synergy coefficient, default penalty amount, and amount of sharing initial real information are larger, the participants are more inclined to keep the agreement. They abide by the rules of the blockchain system, actively share the initial true information, and strengthen the sharing of the initial true information through the incentive and punishment mechanism. &simulation experiments+ sensitivity analysis +the game payoff matrix of the two-party evolutionary game, Jacobian matrix and the sign of the determinant and trace  & some factors affecting the sharing of the initial real information of the AFSC in the blockchain did not considered in the study& NA& China  \\
107 & \textcite{Melesse99} &digital Twins in AFSC & *identify the benefits, types, integration levels, key elements, implementation steps, and challenges of DT in AFSC&AFSC & systematic literature review& DTs can improve efficiency and sustainability in the AFSC by providing visibility, minimizing bottlenecks, planning for contingencies, and improving existing processes and resources &descriptive and content wise analysis   & document availability and coverage&NA& NA  \\
108 & \textcite{Martella101} & ecological Balance of AFSC& assess the environmental sustainability of AFSC  &ISC & empirical approach& overall un-sustainability of the entire supply chain studied &case study , ecological footprint method& lack of detailed data& industrial tomato& Italia  \\
109 & \textcite{Assis103} & Informality in AFSC & factors promoting  informality&LAFSC& qualitative research & Factors that favor informality: family labor and temporary jobs, low education levels, lack of adequate supervision, cultural aspects, or other issues related to necessity and opportunity& interviews, case study, Descending Hierarchical Classification (DHC)& generalization& molluscs  (mussels,oysters and scallops)& Brazil  \\
110 & \textcite{Yontar105} &blockchain and  circular economy of AFSC& research, analyze and prioritize the critical success factors of the use of blockchain technology&CLAFSC & multi criteria approach& blockchain contribute to the circular economy due to “ability to prevent food waste”, “Increased food security”, “Product lifecycle tracking” factors &survery, PESTEL analysis approach, analytic Network Process (ANP),  Multi-Atributive Ideal-Real Comparative Analysis (MAIRCA) & limited expert sample& NA& Turkey  \\
111 & \textcite{Derakhti107} &Challenges and barriers facing the AFSC & identify the challenges and barriers against the application of industry 4.0&AFSC & review of the Literature (SLR)& 15 challenges classified into five major themes: technical, operational, financial, social and infrastructure. The four most important challenges identified are technological architecture, security and privacy, big data management and IoT (internet)-based infrastructure &systematically analyze the literature   & restricted shade of selected items (24) & NA& NA  \\
112 & \textcite{Gholian109} &design of a sustainable AFSC network &simultaneously optimize:total profit, customer satisfaction, production level, water consumption.&AFSC & metaheuristic, convex robust optimization& water consumption is minimal, and a customer satisfaction level of 94.73\%. &stochastic multiobjective
programming model, case study, LP-metric method, hardness of the problems, Keshtel Algorithm, statistical comparison, multi criteria decision-making (MCDM)& generalization& saffron& Iran  \\
113 & \textcite{Yadav110} &IoT in AFSC & examine studies on IoT-based AFSC&AFSC & literature review (network-based cluster analysis & seven identified clusters, with the role of IoT (1:agri-food safety, traceability and sustainability; 2: AFSC sustainability;  3: AFSC performance measurement;  4: AFSC resilience in disruption; 5: AFSC integration and traceability;  6: AFSC transparency and coordination, and  7 identifies the barriers in IoT adoption &  descriptive statistics and network analysis& limited period& NA& NA  \\
114 & \textcite{Dovbischuk111} &logistics service quality & identify attributes of logistics service quality (LSQ) are associated with the superior LSQ &RAFSC, UAFSC  &  expectancy– dis confirmation paradigm in service management research and the related service quality (SERVQUAL) approach&  service quality in agricultural logistics is a five-dimensional construct. Its five dimensions are reliability, digital transformation, corporate image, environmental sustainability, and quality of customer focus.  &interviews, exploratory factor analysis (EFA)   & data collection interrupted by war, generalization, exploratory only approach & NA&Ukraine  \\
115 & \textcite{Peterson112} &impacts of COVID-19 on AFSC businesses & identify changes in production, distribution and consumption following the pandemic&AFSC & statistical approach& impact varies depending on the regions and segments studied &regional survey,  Wilcoxon rank-sum test & results limited to three regions& NA& US  \\
116 & \textcite{Aris114} &impact of technology 4.0 in business performance and  sustainability of AFSC & identify the relationship between the adoption of agricultural revolution 4.0 towards AFSC business performances  and AFSC sustainable agricultural performance encompassing economic, environmental and social dimensions in sustainability&AFSC & literature review& adopting AR 4.0 improves business performance (costs, quality, flexibility and speed) and the sustainability of the SC &conceptual study   & conceptual study, without empirical validation& NA& Malaysia \\
117 & \textcite{Kopishynska116} &application of modern enterprise resource planning systems & explores the possibilities of creating a unified digital information space in a modern cloud based enterprise resource planning (ERP) system to improve the management of all subjects in the territorial community and facilitate the transition to the Industry 4.0 technology landscape&RAFSC & quantitative and qualitative analysis& it is necessary to form a single software ecosystem on an ERP platform that is more flexible to support development &data collection, case study, business process analysis& lack of a secure digital system& NA& Ukraine  \\
118 & \textcite{Montanyà117} &resilience factors of AFSC & extract those factors that have the greatest influence on the degree of resilience of the AFSC&AFSC & qualitative review of THE literature& there are at least three business management strategies, most frequently cited in the scientific literature, that have a positive effect on the resilience of the AFSC. They are: a customer orientation, local distribution and cooperation among agents &integrative review of publications   & specific context of the COVID-19 pandemic& multi-product& multi-country  \\
119 & \textcite{Babu118} & blockchain and Inter-Planetary File System (IPFS) for AFSC traceability   & assessed a system for tracking agricultural products using blockchain's tracing and anti-tampering features&NPAFSC & exploratory approach& the proposed system (Blockchain, IPFS with encryption) improves decentralization and reliability compared to centralized approaches and other classic blockchains &simulation, experimental validation   & limited constraints (non-perishable product)& NA& NA  \\
120 & \textcite{Balezentis119} &viable agri-food supply chains & define Methods and indicators used in the Analysis of the viability of the agricultural supply chain&AFSC & literature review& viability is not based only on short-term orientations and a return to the pre-crisis situation, but also on a long-term evolution towards the "new normal."&art analysis   & lack of empirical data, lack of stakeholder integration& NA& NA  \\
121 & \textcite{Ali120} &risk and resilience in SMEs & define critical risks and strategies to build resilience in SMEs&AFSC & exploratory research & difference in risk and supply chain resilience profiles between developing and developed countries (adoption of digital technologies, training and development and support by national industry bodies)&   interviews, cross comparative analysis& qualitative analysis only& NA& Pakistan, Tanzania, Australia \\
122 & \textcite{Ali121} &digital transformation in AFSC& study impact of digital transformation in Operational Risk Extenuation (supply-demand mismatch, financial and transportation)&AFSC &empirical approach& extenuation of I4Ts help mitigate the impact of critical SC risks on firm performance &survey, statistical analysis, Structural equation model (SEM)   & sample& NA& Australia  \\
123 & \textcite{Philip122} &impact of yield information sharing in AFSC & develop a new eco-label to promote an efficient agricultural method that provides the right balance between yield and environmental impacts through the optimal use of these farm inputs&PAFSC & analytical approach& this model allows to optimize the farm inputs, retail prices, and order quantity so that the stakeholders’ profits are maximized, stakeholders must share yield information when environmentally conscious consumers dominate the market&game-theoretic model   & two echelon, theoretical model, very simplified chain with one farmer, one retailer, one producer& NA& NA  \\
124 & \textcite{Tuni123} &improving environmental sustainability & showcase how assessing the environmental sustainability performance of a multi-tier food SC made up by SMEs can support decisions in order to drive evidence-based green improvements in the SC operations&GAFSC & empirical approach& applying an eco-intensity method to a food chain composed of SMEs allows identifying hotspots and reducing overall environmental impact, while facilitating the implementation of multi-level GSCM over time& eco-intensity  method, longitudinal case study & micro enterprises, limited geographical scope& bread (Patto della Farina)& Italia  \\
125 & \textcite{Ivo124} & roadmap for performance assessment & define a standardized mapping of AFSC to be considered in a PMS, identifying the key stakeholders, material, information, and financial' direct and reverse flows, identify the most critical sustainability-related Performance Areas (PA) (categories) that should be measured in any AF-SSC, and for each PA proposed, detail the main subcategories (characteristics)&GAFSC & analytic approach& the proposed framework to guide the sustainable measurement of different stakeholders CONTAINS, that includes 16 performance areas and 71 subcategories &performance Measurement Systems (PMS), Triple Bottom Line (TBL), interviews, literature review  & until October 2021, only two concepts studied (CE, sustainability )& NA& NA  \\
126 & \textcite{Zhuyue125} &recovery of AFSC network disruption  &analyze disruption propagation and recovery in AFSC by introducing weak ties and evaluating strategies to improve resilience&RAFSC & empirical and case study& combining strengthening existing relationships with establishing new ones enhances resilience &optimization and statistical modeling & generalization& NA& China  \\
127 & \textcite{Stillitano127} & multi-cycle model for measuring the sustainability & propose a multi-cycle approach with expanded assessment boundaries, including co-products, in an attempt to internalize circularity impacts&circular AFSC & mixed-methods approach & customized LC framework proposed can better highlight the environmental and socioeconomic performances of the system of cycles, allowing CE to deliver its promises of sustainability &Life Cycle Assessment (LCA), Environmental Life Cycle Costing (ELCC), Social Life Cycle Assessment (SLCA) in terms of Psychosocial Risk Factor (PRF) impact pathway   & exclusion of distribution, use, and disposal in base LCA& olive oil& Italy  \\
128 & \textcite{Anastasiadis134} &role of traceability & explore the role of traceability in the transition to sustainability driven AFSC&circular AFSC & mixed-methods approach & change at the consumer level drives changes in the supply chain and eventually at the company level &survey , Delphi rounds, Categorical regression   & sample of consumer survey& tomato& Greece  \\
129 & \textcite{Ada144} &Sustainable supplier selection in AFSC & identify the criteria for sustainable supplier selection &DAFSC & mixed-methods approach & performance assessment must rely on the following criteria: material cost, transportation cost, quality assurance, innovation, energy saving, green manufacturing, communication skills, customer satisfaction, productivity, social responsibility, supplier reputation, information technology systems and environmentally friendly technology &literature review, Fuzzy Analytic Network Process (FANP), fuzzy VIKOR   & sample (seven agri-product supplier companies) & NA& Turkey   \\
130 & \textcite{Kumar146} &strategic framework for developing resilience &identify the major impacts of COVID-19 on AFSC and their correlation with different remedial approaches/strategies to improve the resilience of the AFSC&AFSC & mixed-methods approach & supply chain collaboration is the most important strategy, coordination between the stakeholders, information sharing &Quality Function Deployment (QFD) approach, Best-Worst method, Application of Quality Function Deployment (QFD)  &  limited sample& food grains (perishable items)& India  \\
131 & \textcite{Tirkolaee156} &decision-making approach for green supplier selection &  integrate MCDM and robust optimization to determinate the optimum order quantity of each supplier under uncertain conditions &PAFSC &  mixed-methods approach & model validation &AHP- fuzzy TOPSIS,  RGP approach, multi-objective mixed-integer linear programming, robust goal programming  & method based only under economic and GSS criteria, generalization& meat, chicken, schnitzel, rice, fish, tomato paste, lemon juice, tomato, and yogurt& Iran  \\
132 & \textcite{Boenzi129} &environmental impact of fresh products compared to semi-finished products (processing of fresh products by mechanical and thermal methods) in the production process  & evaluate the environmental impact of each step in AFSC&PHL  & empirical approach and life cycle assessment& using fresh products has a lower environmental impact than using semi-finished products, because the perishable nature of the products limits the transport distance of raw materials&impact assessment method, ReCiPe method  & case study, seasonality of products and inter-annual variability of data& Apricot jam& Italy and Europe  \\
133 & \textcite{Assis130} &trust in AFSC & carry out a meta-analysis on the relationships of trust among the stakeholders in agrifood supply chains&AFSC & analysis approach& trust remains a central theme &meta-analysis, bibliometric study  & limited perimeter:  “Agricultural and Biological Sciences” area, "Food Science Technology” , “Agricultural Economics Policy” , “Agriculture Multidisciplinary” , “Agronomy” , “Fisheries”, “Agriculture Dairy Animal Science” , and “Horticulture” & NA& NA  \\
134 & \textcite{Navarro131} &developing hybrid AFSC& identify the stages and the tasks of this short chain, compare them with conventional agri-food chains, approve the hybrid model of AFSC&SASC,  PAFSC & emprical approach& contribution of the hybrid model: the sharing of personnel, infrastructures and services, complementarity in the product range, thus avoiding food waste, knowledge of consumer tastes and needs  &analysi, interviews, study case analysis    & generalization, time-limited assessment (only short term)& fresh fruit and vegetables((peppers, tomatoes, courgettes, cucumbers, aubergines, green beans, watermelons, and
melons)& Spain  \\
135 & \textcite{Prajapati133} &circular economy in AFSC & minimize  collection costs&AFECC & mathematical approach& innovative packaging materials can help to minimize the wastage of food, the cost of collection varies with an increase in the number of nodes& Mixed-Integer Non- Linear Programming (MINLP)   & no multi-echelon durability, simulated data& grains& India  \\
136 & \textcite{Agnusdei135} &research trends  & identify research areas in the field of sustainable AFSC and the role of technological innovation in the transition to a sustainable AFSC&circular AFSC & literature review&  the need for the transition to sustainable AFSC is developing rapidly, digital technologies are scarcely studied within the AFSC sustainability field, post-consumption phase is a challenging issue &Bibliometric, analysis, network analysis and content analysis   & transition to a holistic and multidimensional approach& NA& NA  \\
137 & \textcite{Takavakoglou139} &constructed Wetlands in AFSC & identify the role of constructed wetlands as nature-based solutions in the AFSC of the forthcoming post-COVID-19 era&AFSC &  literature review& numerous applications for constructed wetlands at every stage of AFSC, contribution of constructed wetlands to the strategic goals of the Horizon Europe Strategic Plan  &emergent qualitative analysis (deductive and inductive)   & generalization& NA& Europe  \\
138 & \textcite{Kontopanou141} &triple "A" AFSC model & investigate the impact of agility, adaptability, alignment in AFSC&AFSC & literature review& The triple "A" models promote the increase of the AFSC effectiveness &literature review, analysis and comparison   & selective bibliography study& NA& NA  \\
139 & \textcite{Ali11} &knowledge management practices, risk management culture  & investigate the mechanisms of exposure to supply chain risks&small AFSC &empirical cases& Knowledge application fosters RMC (Risk Management Culture), which enhances SC resilience &quantitative cross-sectional survey, SEM for data analysis& single country  & NA& Australia  \\
140 & \textcite{yadav8} &IoT-based coordination system  &  identify enablers, develop a framework, analyze causal relationships&AFSC &MCDM approach& 'Top Management Support (TMS)' is the main driver for improving coordination &ISM- Fuzzy DEMATEL& expert sample  & NA& India  \\
141 & \textcite{Meemken9} & impacts of sustainability standards & review the existing evidence on sustainability standards&GAFSC &Narrative review methodology& Standards can improve some production practices but are inadequate for achieving widespread food system sustainability and do not significantly advance equity &SLR& limited by the scope, quality, and methodological constraints  & products governed by sustainability standards (e.g., coffee, fruits, vegetables)& NA  \\
142 & \textcite{Joshi10} &AFSC practices & identify the factors influencing the adoption of sustainable agribusiness practices&AFSC &case study& Main barriers: low adoption capacities and absence of cohesive sustainable policy. &semi-structured interviews& sample  & NA& India  \\
143 & \textcite{Kusumowardani12} &Food loss and food waste (FLW) & develop a circular capability framework &AFSC &case study& circular economy adoption is limited; preventative strategies are lacking.  &interviews& sample  & NA& Indonesia  \\
144 & \textcite{Yadav10} & IoT-driven multi-tier system  & analyze cause-effect relationships for an IoT&AFSC &MCDM approach& hierarchical and causal model of enablers  &ISM- Fuzzy DEMATEL& expert sample  & NA& India  \\
145 & \textcite{Gianluca1} &blockchain in AFSC & Examine impact of blockchain technology for traceability, sustainability, and certification in AFSC&AFSC &conceptual approach& blockchain technology can improve supply chain management, profitability, and transparent relationships  &analysis of regulations & Technology cannot replace totally certification processes  & NA& EU/USA  \\
146 & \textcite{Yadav145} &current status and challenges of AFSC and key indicators to measure the performance of AFSC & identify the various challenges in AFSC, review the research contribution in the field of designing AFSC network (AFSCN), investigate performance measurement system of AFSC through various performance indicators&AFSC & literature review& technology adoption and integration in design AFSCN are exciting prospects for future AFSC research &Qualitative content analysis   & consideration only of the concept of sustainability& NA& NA  \\
147 & \textcite{Mehmood147} &drivers and barriers for a circular economy & identify the drivers and barriers for implementing circular economy initiatives &AFSC & literature review& environmental, policy and economy , and financial benefits  are the three top drivers. Institutional , financial , and technological risks  are the top three barriers&Data analysis  & subjectivity in methodology( keywords, inclusion and exclusion criteria)& NA&NA  \\
148 & \textcite{Amentae150} &digitalization in AFSCS& identify the potential contribution of digitalization of the food system &AFSC & literature review& 'Blockchain' was the top hit among keywords, the challenges are infrastructure and cost, knowledge and skill, law and regulations, the nature of the technologies, and the nature of the food system &data analysis   &time scale& NA& NA  \\
149 & \textcite{Rejeb154} &big data in AFSC & identify the contribution of big data to the sustainable development of AFSC & literature review& big data help in crop/plant management, animal management, waste management and traceability &PRISMA    & lack of empirical validation& NA& NA  \\
150 & \textcite{Chiaraluce155} &circular economy in AFSC & identify the trends and future pathways&AFSC & literature review& the topic CE in AFSC is of particular relevance in the scientific community, and the concept CE is continuously evolving & bibliometric analysis   & horizon ends at 2019 and publication at 2021& NA& NA  \\
151 & \textcite{Barbosa160} &research streams on AFSC & identify frequently-used keywords, new and hot research topics and frequently-studied supply chain management practices&AFSC & literature review& frequently used keywords are food supply chain, food waste, sustainability, food safety, SCM, food industry, and food security. New research themes are contract, blockchain, internet of things, resilience, and short food supply chain &Bibliometric analysis   & scope limited to the Web of Science (WoS) database only& NA& NA  \\
152 & \textcite{Galera163} &eco-Innovations and exports interrelationship & analyze the influence between the two variables, both at a business (micro level ) and at a country/region (macro level), and identify the most influential factors ( country of origin and sector of activity)&IAFSC & literature review& there is a positive bidirectional relationship, influenced by several factors, such as social performance, environmental regulation, cooperation strategies, employment level, or business size &systematic literature review    & Methodological limitation( bibliometric approach)& NA& NA  \\
153 & \textcite{Ricardo157} &third party certification of AFSC using smart contracts and blockchain tokens & demonstrate a harvest-level third-party certification (TPC) mechanism for agri-food using Ethereum smart contracts and ERC-1155 tokens to publish tamper-resistant certificates linked to harvests, enabling public access via mobile apps and reducing fraud/green-washing along the supply chain. The approach aims to track origin and quantity of each harvest from farm to consumer, incentivize chain participants, and provide immediate certificate access through a URI on the authority’s site&AFSC & mixed approach& harvest-level TPC via ERC-1155 NFTs with cert links, proven in a PoC on Ethereum to enhance transparency and reduce fraud. &Harvest-level Third Party Certification (TPC) using ERC-1155 NFTs to mint harvest-specific tokens with cert URIs; tokens flow along the supply chain and can be read by consumers via a mobile app; PoC on Ethereum (Rinkeby/mainnet) shows tamper-resistant, transparent provenance and incentives for participant   & generalization& plant agricultural products (e.g., wine grapes) & NA  \\
154 & \textcite{Hajimirzajan158} &large-scale crop planning &  minimize total supply chain costs and guarantee farmers’ minimum income using NGA/RCG coordination&UAFC & mathematical modeling& a large‑scale, climate‑aware crop planning framework that minimizes total supply chain costs   & generalization& potatoes, onions, tomatoes& Iran  \\
155 & \textcite{Firouzeh164} &extended AFSC &  model and analyze an extended sustainable SC by integrating production decisions with consumer behavior, capturing heterogeneity in demand, and evaluating social, environmental, economic, and behavioral performance &organic and conventional AFSC & mixed approach&  including consumer heterogeneity and behavioral dynamics in an extended agro-food supply chain improves understanding of sustainable transitions & agent-based modeling (ABM), discrete event simulation (DES), and system dynamics (SD) & limited empirical data& wine & Australia  \\
156 & \textcite{Rao168} &european private food safety standards & assess how European private food safety standards interact with global AFSC &GAFSC & literature review&the most studied European private food safety standards are owned by retail conglomerates and therefore place the retail sector in a position of influence in the supply chain. These standards influence supply chain structures, market access, and the efficiency of food safety management systems &PRISMA & European focus& NA& EU \\
157 & \textcite{Bottani28} &key performance indicators & determine the map, analyze, and classify KPIs in the AFSC&AFSC & literature review&  KPIs for the Food Supply Chain are systematically reviewed and mapped & bibliometric analysis, content analysis   &  KPIs are derived from only 125 papers accessed on 31 December 2024 and may not be exhaustive& nine  product types (Agricultural products, Meat, Fish, Dairy products, Bread, Animal and plant production, Alcohol-free drinks, Food production, Fruit)& NA  \\
158 & \textcite{Peiyun20} &empirical transformation& identified technology impact in AFSC  to map themes and theories; identify trends and future directions &AFSC  &literature review& five main clusters: behavior intention, sustainability, operation, performance, and risk management. importance of integrating digital technologies to enhance efficiency, transparency, and sustainability in AFSCs &PRISMA, structural equation modeling   & limited scope& NA& NA  \\
159 & \textcite{Kumar126} &sustainable consumption and production & identify circular-based sustainable KPIs for AFSC to achieve sustainable consumption and production, rank KPIs and select best circular alternatives&Circular AFSC & MCDM approach& Economic dimension has highest weight &F-AHP , F-TOPSIS   & generalization& dairy products &India  \\
160 & \textcite{Manikas138} &food security & develop a conceptual framework to assess AFSC  and identify the main global drivers that affect the local food systems &AFSC & conceptual approach& FS indicators to support an early warning system and policy decisions for resilient UAE AFSC &delphi groups, case studies  & generalization& NA& UAE  \\
161 & \textcite{Kusumowardani140} &circular capability  & develop a circular capability framework&UAFC & empirical and case study& proposed framework can enhance the reduction FLW, improve food security, price stability, economic resilience, and environmental benefits &semi-structured interviews , observations   & lack of validation& fruits and vegetables& Indonesia  \\
162 & \textcite{Bager142} &blockchain in AFSC & assess the potential of blockchain technology to promote sustainability in supply chains through increased traceability and transparency and to identify
barriers &IAFC & empirical and case study&  digitalization of supply chain increase efficiency
and reduce costs, disputes, and fraud, while providing more insight end-to-end through product provenance
and chain-of-custody information & semi-structured interviews   & generalization& Coffee& Colombia and Denmark  \\
163 & \textcite{Giovannetti149} &role of cooperation & Assess how Commons and Social Capital through cooperation enhance sustainability and resilience in supply chain&UAFC & empirical and case study& Cooperation and open, rights-based SC improve transparency, pricing fairness, and sustainability &depth case study, interviews  & Single case& Parmigiano Reggiano cheese& Italy  \\
164 & \textcite{Beber152} &strategic actions for internationalization &  investigates strategies for integrating supply chains and how both sides can benefit from this situation also in terms of sustainability&AFSC & empirical and case study& viable sustainable internationalization strategies with emphasis on governance, market entry , differentiation, R\&D training, quality improvements, and supportive policy for win–win outcomes &expert interviews, content analysis, snowball sampling method    & generalization &dairy products& Brazil, Germany  \\
165 & \textcite{Kramer162} &blockchain in AFSC & determine the impact of the three types of blockchain platform in AFSC governance and coordination  &AFSC & empirical and case study& blockchain enhance coordination, information sharing, decision-making, and collective learning benefits &expert discussions   & lack more empirical validations & coffee  and frozen fish & EU   \\
166 & \textcite{Pérez165} &collaborative relationships & analyze collaboration patterns to improve sustainability and performance&PAFSC & empirical and case study& collaboration is more profitable than collaboration with customers &survey, regression   & single regional chain& NA& Spain  \\
167 & \textcite{Ciccullo169} &role of food waste prevention technologies & investigates extensively the range of the available technologies and the detailed objectives of such technologies for food loss and waste prevention&UAFC & empirical and case study&Collaboration with technology providers can be based on continuous technical assistance and consulting for data elaboration and data analysis as well as on full data sharing and co-design, allowing to achieve a different impact on food loss and waste prevention &case analysis   & limited sample&  fruits and vegetables & Italy  \\
168 & \textcite{Zaridis170} &strategy and scale constraints & examine the impact of strategy and scale constraints on AFSC collaboration &AFSC & empirical and case study& collaboration positively affects small and medium-sized enterprises performance & principal components analysis, hierarchical regression, cluster analysis, correlations & limited sample &NA& Greece  \\
169 & \textcite{Cao171} &impact of the green standards at different stages of AFSC & design a cost sharing contract and a buyback contract to coordinate the greening effort decisions &GAFSC & mathematical modeling, theoretical analysis, and numerical simulation & increasing green standards can stimulate the supply chain to increase greening efforts and sales price, which increases costs and decreases product quantities, consequently decreasing the supply chain profits. The greening effort investment by one entity of the supply chain could be affected not only by its own standard, but also by the standard for the other entity &game model (Cobb–Douglas and  demand function)& generic (no distinction between products type)& NA& NA  \\
170 & \textcite{Liu172} &impact of blockchain and big data in coordination of GAFSC & analyze the benefit models of producer and retailer + propose a cost-sharing and revenue-sharing contract to coordinate the AFSCS &GAFSC & mathematical modeling, simulation& when the total investment cost payed by producer and retailer is in a certain range, using ISBD will help chain members gain more benefits. If chain members want to gain more benefits after using ISBD, they should try their best to optimize costs by extracting valuable information &Stackelberg game theory   & SC with one producer and one retailer& NA& China  \\
171 & \textcite{Nattassha173} &circular economy in AFSC & develops a conceptual model of an organic fertilizer producer &GAFSC & soft systems methodology approach& the importance of increasing interaction between the company and the farmers &CATWOE framework   & one echelon& organic fertilizer derived from vegetables and fruit& Indonesia  \\
172 & \textcite{Mehdiyeva174} &barriers and Drivers of the Implementation and Management of GAFSC & analyze the relationship between logistics investments (development of transport and storage capacities) and agricultural performance&GAFSC & quantitative approach& the major driver of the implementation of green supply chain management is the demand for high-quality products and the availability of modern infrastructure. &statistical analysis, documentary analysis   & limited data& NA& Azerbaijan  \\
173 & \textcite{Han175} &home-delivery-oriented agri-food supply chain (HASC) alliance AND THE competitiveness OF small and medium agri-food enterprises & propose a theoretical framework for the HASC alliance and its implementation strategies&LMAFSC & analytic method& the minimum and maximum cost ranges of control strategies at which the alliance can maintain its stability and performance are 5\% and 29\%, respectively, of the total operation cost &literature review, simulation  & generalization& NA& China \\
174 & \textcite{Cui176} &coordinating a green AFSC with revenue-sharing contracts & explore the impact of revenue-sharing contract &GAFSC & mathematical modeling& the revenue-sharing contract is beneficial to increase the greening level, improve both the farmer’s profit and the retailer’s profit.  The effectiveness of the revenue-sharing contract is positively correlated with consumers’ sensitivity to the greening level &Game theory, numerical simulation   & single farmer and single retailer& NA& NA \\
175 & \textcite{Schrobback177} &AFSC assessment & develop a supply chain assessment framework&UAFC & empirical and case study& description matrix and framework& literature review, description and analysis AFSC & limited time series& beef and grain& Australia  \\
176 & \textcite{Fu178} &blockchain in AFSC & study the coupling of management mechanisms between blockchain and the AFSC in the “technology, institution” logic&AFSC &descriptive case studies& blockchain reduces information asymmetry, limits opportunism and is built around organizational structure, trust and contract. &study case  & case-based study&poultry, grain& China  \\
177 & \textcite{Kamble} &data-driven AFSC & propose an application framework for practitioners&AFSC &literature review& Environmental sustainability has the most attention &SLR  & 2000-2017&NA& NA  \\
178 & \textcite{Pancino179} &partnering and the sustainability of AFSC & design of  the process of designing a multi-stakeholder partnership to be implemented in the agri-food supply chain in order to introduce sustainable agricultural practices&MTAFSC & Empirical& it is not possible to design a single standard contract that also includes variation for different locations or division of output between the stakeholders &case study, empirical analysis (documentation, interviews, and participant observations)  & generalization& Barilla products (tomato, durum wheat, and oilseed crops)& Italy  \\
179 & \textcite{Jonkman180} &design AFSC  & integration of seasonality and harvesting decisions, perishability and processing into the design of the AFSC &PAFSC & mathematical modelling& maximize the total gross margin and minimize the global warming potential in CO2 -eq &pareto-efficient frontier, $\varepsilon$-constraint method   &  large differences between time intervals of expected harvest yield& sugar beet& Netherlands  \\
180& \textcite{Cao181} &price formation mechanisms & analyze the price formation mechanisms&upstream
and downstream AFSC & empirical and case study& Price forms at wholesale node and cascades to upstream farmers and downstream retailers/consumers&descriptive statistics   & small sample& fresh and raw vegetables (cucumbers)& China  \\
181 & \textcite{Banasik182} &accounting for uncertainty & propose a multi-objective two-stage stochastic programming model to analyze and evaluate the economic and environmental impacts to account for uncertainty in AFSC&UAFC & mixed approach& applying a 2-stage stochastic programming approach for production planning decisions can lead to improved economic and environmental performance in an AFSC & stochastic programming, deterministic model   & uncertainty limited to yields and demand, single case study&  mushrooms & Netherlands \\
182 & \textcite{Fu183} &Weather risk–reward contract& design a weather risk-reward contract mechanism&AFSC& mathematical approach& this contract can settle the distortion of the sustainable investment level and effectively motivate farmers to participate in the sustainable agricultural practice &Guaranteed Price Mechanism    & lack of social pillar& NA& China  \\
183 & \textcite{Mangla184} &enablers to implement sustainable initiatives &  analyze the key enablers in implementing sustainable initiatives for  AFSC&retail supply chain & empirical case study& three bottom level enablers HAVE the highest independence powers. These enablers are influential enablers. Six enablers were categorized in the cause group and the remaining four in the effect group&literature review, expert consultation, Interpretive Structural Modeling (ISM) - fuzzy Decision Making Trial and Evaluation Laboratory (DEMATEL), MICMAC analysis& the subjective cause, effecT diagram build on the judgments of experts from a particular industry. Further, the identification of the enablers could be challenging& fresh produce (vegetable and fruits), dairy products, and fast-moving consumer goods (FMCG)& India \\
184 & \textcite{Chen185} &blockchain and  Deep Reinforcement Learning in AFSC & design a blockchain-based ASC framework, propose a Deep Reinforcement learning based Supply Chain Management (DR-SCM) method &AFPSC, RSC& experimental, simulation& blockchain ensures reliable product traceability,  DR can achieve higher product products than heuristic and Q-learning methods. &simulation, RL algorithm  & experiences limited to simulated scenarios& corn & NA  \\
185 & \textcite{Stone186} &resilience in AFSC & identify the aspects of resilience&AFSC & literature review& complexity of AFSCs and subsequent exposure to almost constant external interference means that disruptions cannot be seen as a one-off event &SLR  &one aspect& NA& NA  \\
186 & \textcite{Dania187} &collaboration in AFSC & investigate the landscape of extant literature &AFSC & systematic review& 10 key behavioral factors to enable an effective collaboration system for sustainable AFSC :joint Efforts, sharing activities, collaboration value, adaptation, trust, commitment, power, continuous improvement, coordination and stability &resource Dependence Theory (RDT), Theory and Content Analysis , analysis based on frequency& too generic factors, little-studied interactions& NA& NA  \\
187 & \textcite{Luo188} &bibliometric and content analyses & identify the homogeneous areas in AFSC and to assess the movement and interactions within and between fields &AFSC & bibliometric and content analyses& 6 clusters are identified : short/alternative food supply chains, ASC Sustainability, modeling ASC traceability, risk management, and optimization, global ASCs, ASC transparency and traceability, ASC relationships/vertical coordination/networks& systematic literature review combined with bibliometric and content analyses & subjectivity of article selection& NA& NA  \\
188 & \textcite{Pappa189} &electronic traceability systems (ETsystems) &investigate the factors influencing the acceptance and use of ETsystems in AFSC&GAFC & mixed approach&‘Perceived Control’ and most importantly, the ‘perceived costs’ over the installation and operation of the ETsystem, is the most important factor with the strongest direct effect influencing the intention to install and operate such a system & technology Acceptance Model  (TAM2), theory of Planned Behavior (TPB), PLS-SEM analysis & generalization & dairy products & Greece  \\
189 & \textcite{Zhao190} &theory of AFSC resilience & identify AFSC resilience factors (enablers, model and their interrelationship) and develop a comprehensive resilience framework &AFSC& mixed approach&  leadership and working-team stability are top driving factors; traceability and information sharing are critical but largely dependent on others  &structural modeling total interpretative and cross-impact multiplication matrix   & limited interview coverage & NA& United Kingdom, France, Spain, Italy, Argentina  \\
190 & \textcite{Allaoui191} &sustainable AFSC design& optimizing the three pillars of sustainability&AFSC & mixed approach& two-stage hybrid approach yields a Pareto frontier showing trade-offs between cost, emissions, water use, and jobs & AHP, OWA for partner efficiency, multi-objective MILP, Pareto front via $\varepsilon$-constraint  &  data limited to case study& NA& NA  \\
191 & \textcite{Stefanella192} &blockchain and AFSC performance &  assess impact of blockchain on the AFSC performance& PAFSC & empirical and case study& blockchain improve information accessibility, availability and sharing. However, the flexibility and responsiveness of operations within the AFSC do not impacted &case study, thematic analysis   & generalization, limited sample& Poultry meat, lemons, and oranges & Europe  \\
192 & \textcite{Qifan193} &introducing of online channel in AFSC & investigate the practicality of introducing online channels in the AFSC&GAFSC & mathematical optimization&the green food brand effect of online channels is does not necessarily improve with scale, and the initial freshness has a positive relationship to the profit, demand, and price of farmer cooperatives and online retailers &game theory   & constant price during product shelf life& NA& NA  \\
193& \textcite{Chen194} &bio-refinery technologies for unused biomass & impact of bio-refinery technologies for unused biomass in sustainability of AFSC &AFPSC & experimental approach& improving agricultural production efficiency is not enough to achieve sustainable global food security, developing bio-refinery technologies for unused biomass in the agri-food supply chain could be the crucial factor &technique practical experimental & laboratory tests& green tea residue& china  \\
194 & \textcite{Garofalo195} &environmental sustainability of AFSC & identify the impacts related to greenhouse gas emissions &AFSC &  International Reference Life Cycle Data System (ILCD) &waste produced during the processing phase was the main contributor to the total environmental impact, followed by the packaging and cropping phase &life cycle assessment methodology & partial sectoral study& whole-peeled tomato& Italy  \\
195 & \textcite{Tasca199} &environmental sustainability of AFSC & study the combination of organic farming and alternative distribution (local and large-scale) and their impact to reduce the environmental footprint compared to the conventional model &RAFSC & empiric& no agricultural techniques have a better overall environmental profile, distribution of raw organic product loose  is generally the most advantageous in the majority of impact categories&life cycle assessment LCA & limited perimeter& endive& Italy  \\
196 & \textcite{Akhtar200} & adaptive leadership and sustainable AFSC& study the link between adaptive leadership and sustainability &IAFSC & Descriptive statistics& data-driven and adaptive leadership is a key determinant for sustainability &structural equation modeling(SEM) & survey research& dairy,meat,vegetables and fruits&from/to China,India and New Zealand and another country\\
197 & \textcite{Aggarwal201} &collaboration in AFSC & study the impact of collaboration on AFSC efficiency &AFSC & empirical approach, qualitative methodology & collaboration brings benefits to buyers and suppliers as well as the application of sustainable best practices in the agi-food industry (reduction of waste, price vocation, etc) &case Study, interview, grounded theory methodology & single case study, generalization& rize& India  \\
198 & \textcite{Accorsi202} &ecosystem carbon balance in AFSC & develop A linear programming model to optimize agriculture, and logistics costs  and minimize the carbon emissions  &AFSC & mathematical modeling& planning focused solely on profitability leads to imbalances in the sustainability of agri-food chains, hence the importance of integrating environmental criteria &linear programming (LP), study case & scale restriction & potato& Bologna  \\
199 & \textcite{Galal203} &uncertainty of demands in AFSC & study the effect of changing the order quantity and lead time in AFSC on economic and environmental performance measures (stock levels, associated costs and CO2 emissions)&PAFSC & analytical simulation approach& reducing order quantities can reduce costs and emissions; in addition, it is important to integrate the emissions because costs IS not sufficient for chain performance  &discrete event simulation model, case study& data reliability (developing countries)& oranges& from Egypt to  Netherlands  \\
200 & \textcite{Akhtar204} &leadership and coordination between AFSC ACTORS & impact of participative leadership  and  directive leadership in AFSC coordination&IAFSC & empirical approach& participative leadership approach is highly significantly  associated with the effectiveness of AFSC coordination unlike the directive leadership &survey, covariance-based structural equation modelling (CBSEM), exploratory factor analysis, reliability, and validity tests & generalization& dairy, meat, apples, onions, and wine& New Zealand-UK  \\
201 & \textcite{Rota205} &AFSC governance structures and sustainability & determine  the impact of  processes oriented sustainability  on the supply chain governance structures  &AFSC & empirical approach& the integration of sustainability dimensions into the governance structures of agri-food chains, strengthening technical and logistical efficiency and the improvement of food quality&adoption of experts’ evaluations & complexity, generalization& soy and beef& Latin America  \\
202 & \textcite{Folinas206} &lean thinking practices in green AFSC & propose a systematic approach to determine waste (by measuring the time without added value in production and logistics processes, as well as the water and energy consumed)  &GAFSC & systematic approach&  integration of Lean techniques, VSM, and ERP modules  improves the performance of the AFSC &study case, value stream mapping (VSM) & limited generalization& corn for animal feed&  Greece   \\
203 & \textcite{Tsolakis207} &decision making in AFSC & recognize the hierarchy of the decision-making process for the design and management of AFSCs (Strategic, tactical and operational levels)&AFSC & literature review& decision-making processes in the agri-food supply chain are dynamic and stochastic in nature; need to develop real-time optimization tools to solve dynamic and stochastic problems in the AFSC &synthesis of the literature   & static logistics chain& NA& EU  \\
204 & \textcite{Verdouw210} &TIC in AFSC &design a platform for smart agri-food logistic information systems &AFSC & user-oriented approach& TIC facilitates real time virtualization, logistics connectivity, and logistics intelligence of AFSC &study case (needs, specifications, technical design)   & generalization& fresh fruit and vegetables (FFV) industry and the Plants and Flowers (PF) industry& EU  \\
205 & \textcite{Boudahri211} &configuration AFSC network & develop a clustering-based location-routing model &PSL & mathematical modeling approach& minimize the total costs of the AFSC, determine the location of the slaughterhouses, and determine the distribution routes to the customers &Branch and Bound   &one echelon of AFSC& chicken and turkey& Algeria  \\
206 & \textcite{van213} &integration of product quality in TCAFSC & develop a preliminary diagnostic instrument, which can be used to assess supply chain networks on QCL(quality controlled logistics) and identify opportunities for improvement&TCAFSC & exploratory case study& using time-depend quality information helps to optimize product quality and availability in markets and minimize shrinkage &literature review, study case   & limited sectoral scope& tomato& Netherlands  \\
207 & \textcite{Fischer214} &contractual choice and sustainable relationships  &identify factors influencing the choice of contract type and which lead to sustainable inter-enterprise relationships&DTAFPSC & empirical approach& market, industry, and enterprise-specific characteristics influence the choice of contract type &survey, statistical data analysis (binary logit model + structural equation model SEM))  & limited sectoral scope& pig meat, beef, and cereals& Germany, UK, Spain, Ireland, Finland, and Poland  \\
208 & \textcite{aramyan217} & performance measurement & develop a generic framework and evaluate it & CBSC & conceptual framework, case study & identifies 4 main categories: efficiency, responsiveness, food quality, flexibility &conceptual framework analysis &absence of empirical validation and integration of multi-criteria indicators & tomato & Netherlands \& Germany \\
209 & \textcite{Sharma2020} & Machine learning (ML) applications in sustainable AFSC& Identify the current state and future potential of ML applications in AFSC & AFSC& literature review&Artificial neural networks (ANN) are the most frequently used ML algorithm&Systematic literature review + descriptive statistic & 2000‑2019& NA &  NA\\
210 & \textcite{Kirti2023} & blockchain in AFSC&identify the factors that affect blockchain technology adoption in the AFSC& AFSC & empirical approach& blockchain technology adoption  has a positive effect on sustainable supply chain performance  &SEM&sample & fresh fruits, vegetables, beverage, dairy & India\\
211 & \textcite{Shashank2021} & circular economy adoption barriers  &  identify ranking of barriers& AFSC & Multi-criteria approach& Most critical barriers are lack of government support and incentives (LGS) and lack of effective policy and protocol (LPP) &Delphi, ISM, and ANP&Generalization & pulp, milk, jute , rice, sugar cane, wheat, cotton, vegetables, and fruit & India \\
212 & \textcite{Sachin} &  Blockchain technology in AFSC& identify the enablers that encourage blockchain adoption in AFSC  & AFSC&MCDM approach& Most significant enabler is traceability&Delphi, ISM‑MICMAC &Generalization & NA & India \\
\end{longtable}
\endgroup
\end{small}
\end{landscape}
\printbibliography
\end{document}
